% !TeX root = main.tex
\chapter{Аналіз методів виявлення аномальних кодів у фото- та відеооб’єктах вебресурсів}
\label{chap:analysis}

Фото- та відеооб’єкти надходять до вебсистеми через штатні канали приймання вмісту, після чого їх перевіряють, декодують, перетворюють, зберігають і передають користувачеві. Такий маршрут створює особливу умову прихованого впливу. Об’єкт може відповідати обмеженням формату й не мати помітних зорових спотворень, але зберігати дані, які інший компонент зчитує, відновлює або використовує під час оброблення. Тому успішне декодування характеризує лише формальну придатність носія й не є достатнім свідченням його безпечності.

Причина цього обмеження полягає в багаторівневій будові мультимедійних форматів. Пікселі та декодовані кадри утворюють зорове представлення; коефіцієнти перетворень і локальні залишки формують частотне та вейвлетне представлення; сегменти, блоки, маркери, потоки й таблиці визначають структуру контейнера; службові поля можуть зчитуватися незалежно від зорового вмісту. Для відео додаткове значення мають часові зв’язки, параметри прогнозування, вектори руху й рішення кодера. Огляди стеганографічного аналізу показують, що детектор, розроблений для окремого носія або способу вбудовування, не набуває автоматично здатності контролювати всі ці представлення \cite{kheddar2024review,apau2024review}.

Наявні дослідження відповідно зосереджуються на різних джерелах свідчень. Для зображень аналізують слабкі просторові й частотні зміни, для відео додатково враховують часову організацію та ознаки стисненого потоку \cite{muralidharan2022race,bouzegza2022video}. Засоби ідентифікації байтових фрагментів за вмістом і виявлення багатоформатних об’єктів перевіряють внутрішню будову незалежно від форматного розширення \cite{mittal2021fifty,koch2022polyglot}. Поведінкове спостереження, своєю чергою, може зафіксувати реакцію компонента, якої немає у статичному поданні. Ці напрями розв’язують взаємопов’язані, але не тотожні задачі та використовують непорівнювані за змістом ознаки.

Для систематизації доступного корпусу джерела 2021--2026~рр. зіставлено за п’ятьма критеріями: видом носія й формату, рівнем представлення, механізмом формування ознак, типом результату та обмеженнями експериментального протоколу. До аналізу включено праці, що описують принаймні одну з чотирьох функцій: приховане розміщення даних; їх виявлення або локалізацію; ідентифікацію байтових фрагментів чи форматної неоднозначності; спостереження за проявом під час оброблення. Технічні стандарти й офіційні записи про уразливості використано лише для уточнення правил формату та меж оброблення, а не як джерела класифікаційних показників.

Огляд виконано як аналітично-систематизувальний. Протокол формально систематичного огляду з повним журналом баз, пошукових виразів, правил вилучення та усунення дублювань не формувався; тому зіставлення встановлює охоплення й взаємні обмеження включених праць, але не доводить вичерпності корпусу. Оскільки набори даних, способи приховування й протоколи перевірки відрізняються, опубліковані числові результати не агреговано в спільний рейтинг.

Аномальним кодом у роботі названо приховану послідовність, значення або структурну вставку, що не відповідає заявленому призначенню носія та створює умови для несанкціонованої дії; розгорнуте означення й розмежування з поняттями аномалії, прихованої загрози та шкідливого коду подано в підрозділі~\ref{sec:hidden_code_placement}. Предметну область обмежено растровими фото- та відеооб’єктами поширених у вебресурсах форматів --- JPEG, PNG і WebP для зображень та контейнерів родини MP4 і WebM для відео. Векторні (SVG), анімовані (GIF, APNG), а також новіші формати зображень (AVIF, HEIC) і додаткові відеокодеки (H.264/AVC, VP9, AV1) визначають напрями подальшого розширення, але залишаються поза межами перевіреного в роботі обсягу.

У розділі визначено умови використання мультимедійних об’єктів як носіїв прихованого впливу, систематизовано способи розміщення аномальних кодів, зіставлено відомі групи методів і встановлено вимоги до подальшого дослідження. Аналіз побудовано від шляху об’єкта у вебсистемі до наукової прогалини, яка потребує узгодженого опрацювання структурних, спектрально-зорових і подієво-поведінкових свідчень.

\section{Аналіз використання мультимедійних об’єктів вебресурсів як засобів прихованого впливу}
\label{sec:multimedia_threat_carrier}

Мультимедійний об’єкт бере участь у послідовності автоматизованих операцій і на кожному етапі набуває іншого доступного представлення. Модуль приймання зіставляє заявлений тип, розширення, сигнатуру й обмеження розміру. Декодер застосовує синтаксичні правила конкретного формату. Засіб керування вмістом може читати службові поля, а клієнтське середовище використовує результат під час побудови сторінки. Розбіжність між цими інтерпретаціями дає змогу одному компоненту визнати об’єкт допустимим, тоді як інший компонент отримає доступ до прихованої частини даних.

У межах цієї роботи фото- або відеооб’єкт розглядається як носій, а не як самостійна виконувана програма. Прихована послідовність проходить штатним каналом до компонента, здатного її прочитати або відновити. Наслідок залежить від місця розміщення, представлення даних, сценарію оброблення та прав компонента. Зображення, яке коректно відображається, може передавати команду пов’язаному процесу, містити додаткову структуру після логічного завершення або зберігати фрагмент, який відновлюється зовнішньою логікою. Дослідження приховування шкідливого коду в зображеннях підтверджує відмінність між звичайним секретним повідомленням і програмним навантаженням, для якого потрібне окреме виявлення \cite{kumagai2022malwareimage}.

Для оцінювання такого носія доцільно розмежовувати наявність прихованої послідовності, доступність цієї послідовності певному компоненту та можливість несанкціонованої дії. Сам факт нетипових даних характеризує властивість об’єкта, але ще не встановлює їхнє призначення. Доступність виникає тоді, коли на маршруті оброблення є засіб, здатний інтерпретувати відповідне подання: байтову область, службове поле, піксельні значення або параметри стисненого потоку. Можливість дії додатково залежить від сценарію, прав і зовнішньої логіки, яка використовує відновлені дані. Саме зв’язок носія з такою логікою відрізняє потенційно прихований вміст від реалізованого каналу впливу \cite{iglesias2023driveby}. Водночас відсутність спостережуваної дії під час одного відкриття не спростовує доступності прихованої послідовності за іншого сценарію. Отже, первинний аналіз повинен не приписувати знайденому фрагменту наперед визначену семантику, а встановити, де він розміщений, який компонент може його прочитати та на якому відрізку маршруту це відбувається.

Шлях такого носія від зовнішнього джерела до можливого прояву прихованого впливу наведено на рис.~\ref{fig:multimedia_threat_path}.

\begin{figure}[H]
\centering
\begin{tikzpicture}[
  x=1cm,
  y=1cm,
  figure block/.style={dissertation block,minimum height=1.55cm}
]
\node[figure block,text width=2.75cm] (src) at (-6,2) {Зовнішнє\\джерело};
\node[figure block,text width=2.90cm] (obj) at (-2,2) {Цифрове зображення\\або відеозапис};
\node[figure block,text width=3.25cm] (recv) at (2,2) {Приймання та\\форматна перевірка};
\node[figure block,text width=3.70cm] (proc) at (6,2) {Контрольоване розкодування\\та перетворення};
\node[figure block,text width=3.60cm] (orig) at (2,-0.8) {Оригінальний контейнер\\і службові дані};
\node[figure block,text width=3.25cm] (store) at (6,-0.8) {Зберігання та\\публікація};
\node[figure block,text width=4.20cm] (effect) at (6,-3.4) {Прояв\\прихованого впливу\\або формування ознак};

\coordinate (effectroute) at ($(proc.east)+(0.55,0)$);

\draw[dissertation arrow] (src.east) -- (obj.west);
\draw[dissertation arrow] (obj.east) -- (recv.west);
\draw[dissertation arrow] (recv.east) -- (proc.west);
\draw[dissertation arrow] (recv.south) -- (orig.north);
\draw[dissertation arrow] (orig.east) -- (store.west);
\draw[dissertation arrow] (proc.south) -- (store.north);
\draw[dissertation arrow] (store.south) -- (effect.north);
\draw[dissertation arrow] (proc.east) -- (effectroute)
  -- (effectroute |- effect.east) -- (effect.east);
\end{tikzpicture}
\caption{Маршрут цифрового зображення або відеозапису у вебсистемі та точки формування ознак прихованого впливу}
\label{fig:multimedia_threat_path}
\end{figure}


Рисунок відокремлює місце розміщення даних від місця їхнього прояву. Піксельну або частотну послідовність може зчитувати компонент, для якого зображення є каналом передавання, а структурну вставку може ігнорувати засіб перегляду, але обробляти інший інтерпретатор. Отже, аналіз має встановлювати не тільки область приховування, а й шлях, на якому ця область стає доступною.

Для фотооб’єктів кілька представлень виникають у межах одного формату. JPEG поєднує маркерну структуру, блокове дискретне косинусне перетворення (DCT) та квантування, тому його можна досліджувати як матрицю пікселів, сукупність частотних блоків, послідовність сегментів і набір службових полів. PNG зберігає точні значення після безвтратного декодування та використовує впорядковані блоки з контрольними сумами. WebP має контейнерну організацію формату RIFF і кілька режимів кодування; спеціалізований протокол приховування для WebP показує залежність ознак від внутрішньої будови формату \cite{koptyra2023webp}. Перелічені властивості самі по собі не є ознаками загрози, але визначають допустимі варіації, з якими потрібно зіставляти спостережувані зміни.

У відео розбіжність між зоровим результатом і стисненим поданням є ще суттєвішою. Контейнер організовує потоки, часові мітки, індекси, параметри синхронізації та службові таблиці, а кодек використовує внутрішньокадрове й міжкадрове прогнозування, перетворення залишків і поділ блоків. Прихована послідовність може міститися в окремих кадрах, бути розподіленою в часі або змінювати параметри стисненого потоку. Тому покадрове декодування не зберігає повної сукупності початкових ознак. Огляд відеостеганалізу підтверджує залежність відповідних методів від стандарту кодування та способу вбудовування \cite{bouzegza2022video}.

Формальна перевірка типу, сигнатури, розміру, геометричних параметрів і можливості декодування відсіює частину некоректних об’єктів, але не охоплює всі канали. Коректна початкова сигнатура не виключає даних після завершального маркера, а синтаксична правильність основних блоків не усуває можливості іншої форматної інтерпретації. Випробування засобів ідентифікації типу на багатоформатних об’єктах показало обмеженість перевірки лише за зовнішніми атрибутами \cite{koch2022polyglot}.

Розбіжність інтерпретацій посилюється тим, що етапи перевірки й використання часто виконують різні програмні компоненти. Модуль приймання може зіставити розширення та початкову сигнатуру, засіб декодування прочитати лише підтримувані сегменти, засіб вилучення службових даних перейти до вкладеного подання, а клієнтський компонент застосувати власні правила до отриманого результату. Об’єкт, відхилення якого проігнорував один засіб, через це не обов’язково буде нейтральним для іншого. Виявлення багатоформатності підтверджує саму можливість різних трактувань байтової послідовності, тоді як відомості про уразливість декодера WebP показують безпекове значення межі, на якій дані переходять до конкретної реалізації \cite{koch2022polyglot,nvd2023webp}. Ці приклади не дають підстав вважати кожну різницю між засобами загрозою. Вони обґрунтовують іншу вимогу: результат зовнішньої перевірки потрібно співвідносити з внутрішньою граматикою та з фактичним шляхом подальшого використання.

Зорова правдоподібність також не є достатньою. Користувач сприймає предметний зміст, колір і послідовність кадрів, тоді як програмний компонент читає довжини сегментів, таблиці, службові значення, коефіцієнти та байтові області. Мале або розподілене навантаження може не створювати помітного спотворення, особливо в текстурних ділянках. Відсутність зорових змін тому характеризує непомітність, але не безпечність.

Службові дані формують окремий канал. EXIF, XMP, IPTC, коментарі, профілі кольору та допоміжні поля описують походження й параметри оброблення; склад полів IPTC визначено технічним стандартом \cite{iptc2024metadata}. Вебсистема може зберігати ці значення, передавати їх іншому сервісу або використовувати для формування похідного об’єкта. Нетиповий рядок, велике поле чи зовнішнє посилання потребують перевірки синтаксису, призначення та фактичного використання, оскільки сам розмір або рідкісність поля не доводять наявності загрози.

Зв’язок носія з подальшою дією у середовищі вебсистеми досліджено на прикладі коду, прихованого у зображеннях стеганографічними й багатоформатними засобами \cite{iglesias2023driveby}. Для цієї роботи зазначене джерело підтверджує можливість такого зв’язку, але його числові результати не переносяться на інший набір даних. Практичний наслідок полягає в необхідності прив’язати статичні ознаки та події оброблення до того самого об’єкта.

У життєвому циклі носія змінюється доступність свідчень. Під час приймання наявний повний байтовий потік, але ще немає нормалізованого зорового подання. Передавання декодеру відкриває пікселі або кадри та водночас створює ризик втрати невідомих службових конструкцій. Уразливість декодування WebP CVE-2023-4863 ілюструє, що формально допустимий мультимедійний об’єкт може становити небезпеку саме на цій межі оброблення \cite{nvd2023webp}. Масштабування й перекодування формують похідні копії, зберігання відокремлює об’єкт від початкового контексту, а публікація передає його клієнтським компонентам. Контроль в одній точці не забезпечує простежуваності на всьому маршруті.

Початкове байтове подання необхідно зберегти до нормалізації, оскільки лише в ньому доступні дані після логічного завершення, початковий порядок блоків, вкладені сигнатури та суперечності між типом і граматикою. Декодоване подання, навпаки, потрібне для локальних залишків, частотних і вейвлетних залежностей. Мініатюра, попередній перегляд і перекодований фрагмент повинні зберігати зв’язок із джерелом через ідентифікатор, контрольний відбиток і журнал перетворень. Інакше похідні копії можуть бути хибно витлумачені як незалежні спостереження під час оцінювання або експлуатації.

За характером впливу на свідчення перетворення можна розглядати як незмінне перенесення, безвтратне приведення до узгодженого подання та формування похідного подання з утратою частини початкових даних. Незмінне перенесення зберігає байтову структуру, але може відокремити об’єкт від транспортного контексту. Безвтратне декодування та повторне збереження переважно підтримують піксельні значення, проте можуть змінити порядок блоків або склад службових полів. Масштабування, зміна колірного подання й кодування з утратами формують нові коефіцієнти та здатні послабити просторовий або частотний слід. Для відео перекодування додатково змінює рішення прогнозування, вектори руху та часову організацію потоку \cite{bouzegza2022video}. Дослідження протоколів WebP, витоку початкових даних через ескізи та трас походження підтверджують, що значення ознаки потрібно тлумачити разом з історією формування конкретної копії \cite{koptyra2023webp,benes2023thumbnails,kiltz2024trace}. Жоден клас перетворення не гарантує автоматичного усунення загрози: зруйнований прихований вміст може залишити статистичний слід, а очищене зорове подання не відображає структурних властивостей початкового об’єкта. Тому журнал походження має зв’язувати копії, але не підміняти окреме оцінювання доступних у кожній із них свідчень.

До публікації доцільно виконувати перевірки, потрібні для допуску об’єкта, а суперечливий результат спрямовувати на поглиблений аналіз у контрольованому середовищі. Доступність ознак та відповідні безпекові дії на основних етапах узагальнено в табл.~\ref{tab:lifecycle_evidence}.

\begin{table}[H]
\caption{Доступність ознак прихованого впливу на етапах життєвого циклу мультимедійного об’єкта}
\label{tab:lifecycle_evidence}
\centering
\scriptsize
\begin{tabularx}{\textwidth}{>{\raggedright\arraybackslash}p{0.18\textwidth}>{\raggedright\arraybackslash}p{0.22\textwidth}YYY}
\toprule
\textbf{Етап} & \textbf{Представлення, доступне для контролю} & \textbf{Ознаки, що зберігаються найкраще} & \textbf{Ознаки, які можуть бути втрачені} & \textbf{Безпекова дія} \\
\midrule
Приймання початкового об’єкта & Повний байтовий потік, заявлений тип, транспортні атрибути & Сигнатури, маркери, блоки, службові поля, додаткові області & Ще не сформовано декодовані спектрально-зорові ознаки & Збереження відбитка, структурний розбір, перевірка обмежень \\
Контрольоване декодування & Початковий об’єкт і нормалізоване піксельне або кадрове подання & Локальні залишки, частотні й вейвлетні залежності, помилки декодера & Частина рідкісних службових конструкцій може бути проігнорована декодером & Ізоляція процесу, вилучення ознак, журналювання подій \\
Масштабування або перекодування & Похідне зорове подання та журнал перетворення & Семантичний зміст, частина стійких багатомасштабних відхилень & Молодші біти, початкові коефіцієнти, службові та додаткові області & Пов’язування похідної копії з початковим об’єктом, повторне оцінювання \\
Зберігання і кешування & Початковий об’єкт або похідна копія, службові дані сховища, права доступу & Раніше вилучені ознаки та рішення за умови незмінності & Контекст походження за відсутності журналу; ознаки початкового об’єкта, якщо збережено лише копію & Карантин, контроль цілісності, обмеження доступу до рішення \\
Публікація і споживання & Відповідь вебсервера, клієнтське подання, події взаємодії & Поведінкові та контекстні прояви конкретного сценарію & Внутрішня структура може бути недоступною клієнтському контролю & Застосування політики, блокування або безпечна реконструкція \\
\bottomrule
\end{tabularx}
\end{table}

Жодний етап не зберігає всі свідчення одночасно. Початковий об’єкт потрібний для перевірки структурної цілісності, контрольоване декодування відкриває спектрально-зорові й процесні прояви, а журнал перетворень підтримує причинний зв’язок між ними. Отже, багаторівневий аналіз зумовлений не лише різними способами приховування, а й незворотною зміною доступної інформації під час штатної роботи вебресурсу.

Предметні рівні, які потрібно розрізняти до подальшої формалізації моделі мультимодального представлення прихованої загрози для фото- та відеооб’єктів вебресурсів, наведено в табл.~\ref{tab:object_levels}.

\begin{table}[H]
\caption{Рівні представлення мультимедійного об’єкта з погляду прихованого впливу}
\label{tab:object_levels}
\centering
\small
\begin{tabularx}{\textwidth}{>{\raggedright\arraybackslash}p{0.19\textwidth}YY}
\toprule
\textbf{Рівень представлення} & \textbf{Дані, що підлягають аналізу} & \textbf{Можливий прояв прихованого впливу} \\
\midrule
Зоровий & Піксельні значення, канали кольору, прозорість, локальні залишки, декодовані кадри & Слабкі просторові зміни, розподілене вбудовування, локальна невідповідність шумового профілю \\
Частотний і вейвлетний & DCT-коефіцієнти, вейвлетні піддіапазони, високочастотні залишки, багатомасштабні залежності & Зміни, пристосовані до стиснення або текстури й непомітні у звичайному зоровому поданні \\
Структурний & Заголовки, сегменти, блоки, маркери завершення, таблиці, атоми контейнера, байтові фрагменти & Додаткові області, суперечливі довжини, друга форматна інтерпретація, нетиповий порядок елементів \\
Службовий & EXIF, XMP, IPTC, коментарі, профілі, допоміжні поля формату & Закодовані команди, посилання, надмірні або семантично невідповідні значення \\
Часовий & Послідовність кадрів, ключові кадри, параметри прогнозування, вектори руху, режими поділу блоків & Розподілення даних між кадрами або зміна параметрів стисненого відеопотоку \\
Поведінковий & Операції декодування й перетворення, звернення до мережі та сховища, створення процесів, події пам’яті & Прояв прихованої дії під час контрольованого оброблення або взаємодії з компонентом вебсистеми \\
\bottomrule
\end{tabularx}
\end{table}

Властивість рівня представлення, виявлене відхилення та висновок про приховану загрозу не є тотожними поняттями. Виміряна ентропія, довжина поля, частотний залишок або системна подія становлять первинну ознаку. Її невідповідність очікуваному профілю формує підставу вважати спостереження аномальним, але така невідповідність може виникати через допустиме розширення формату, особливість пристрою, редакторське перетворення чи штатну роботу компонента. Висновок про загрозу потребує встановлення того, що відхилення суперечить призначенню носія та узгоджується з іншими свідченнями можливого впливу. Це розмежування запобігає двом крайнім тлумаченням: автоматичному блокуванню кожного рідкісного об’єкта та допуску формально коректного носія за відсутності зорових спотворень. Отже, рівні в табл.~\ref{tab:object_levels} визначають джерела спостережень, а не готові класи рішення. Їх потрібно аналізувати з урахуванням умов отримання й взаємного підтвердження.

Статичні й процесні дані виникають у різний час і мають різні умови достовірності. Структурні суперечності доцільно встановлювати до перетворення контейнера. Спектрально-зорові ознаки потребують піксельного або кадрового подання. Поведінкові прояви залежать від сценарію, версії засобу, прав доступу й оточення. Їх не можна механічно звести до одного початкового фільтра без збереження походження та умов отримання.

Таким чином, фото- та відеооб’єкт стає засобом прихованого впливу завдяки поєднанню штатного каналу надходження, зовнішньої правдоподібності та неоднозначності внутрішньої інтерпретації. Подальший аналіз має простежувати носій від приймання до компонента, здатного прочитати або перетворити приховані дані. Для визначення потрібних ознак способи розміщення аномальних кодів необхідно систематизувати за рівнями представлення.

\section{Аналіз способів прихованого розміщення аномальних кодів у фото- та відеооб’єктах}
\label{sec:hidden_code_placement}

Стеганографічне перетворення не можна автоматично ототожнювати зі шкідливою діяльністю. Однакові математичні принципи застосовують для конфіденційного передавання, цифрового маркування, захисту авторства й прихованого передавання команд. У цій роботі аномальним кодом вважається прихована послідовність, значення або структурна вставка, що не відповідає заявленому призначенню об’єкта та створює умови для несанкціонованої дії, відновлення керувальних даних чи небезпечної взаємодії з вебсистемою. Спосіб вбудовування описує канал, тоді як висновок про загрозу потребує оцінювання вмісту, структури й контексту оброблення.

Через це назва алгоритму не є достатньою основою класифікації. Молодші біти можуть містити водяний знак або команду, службове поле може бути штатним описом або навмисно сформованою областю, а зміна вектора руху може належати законному маркуванню чи прихованому каналу. Доцільно розрізняти способи за рівнем представлення, який змінюється, та за етапом можливого зчитування. Таке групування відповідає оглядам, у яких просторові, частотні й відеоспецифічні способи розглядаються окремо через різну природу їхніх слідів \cite{muralidharan2022race,bouzegza2022video}.

Механізм просторового розміщення полягає у зміні піксельних значень або прозорості. За рівномірного вбудовування у молодші значущі біти, табличного кодування та поєднання з шифруванням спостерігаються бітові й локальні піксельні відхилення \cite{panigrahi2025lsb,hapsari2023lsb,huang2025hybridcoding,ali2024videolsb,sakshi2022lsb,alanzy2023lsb}. Адаптація до країв, текстури, локальних двійкових шаблонів, апаратних обмежень або кількох рівнів переносить слід у залежності між сусідніми елементами й шумовим профілем \cite{hassaballah2021iiot,sultana2023edge,hameed2024medical,harb2023hardware,abuali2024multilevel,wozniak2025twophase,alexan2024stegocrypt}. Отже, межею одного просторового детектора є правило вибору позицій: повнота для рівномірної зміни не переноситься на адаптивне вбудовування без окремої перевірки. Для власного рішення це зумовлює потребу зберігати локальні залишки й просторові координати, а також історію перетворень конкретного об’єкта.

Спостережуваність просторового сліду визначається подальшим маршрутом. Безвтратне копіювання може залишити молодші розряди незмінними, тоді як масштабування, корекція кольору та повторне кодування з утратами здатні зруйнувати приховану послідовність, але залишити змінені локальні розподіли або нетипові залишки. Через спільне походження початковий і похідний об’єкти не є незалежними прикладами. Тому метод має реєструвати перетворення й не розподіляти споріднені копії між незалежними частинами оцінювання.

Механізм частотного розміщення змінює коефіцієнти перетворення зорових даних. Для JPEG спостережуваними є розподіли квантованих DCT-коефіцієнтів, параметри квантування й міжблокові залежності, а характер сліду залежить від правила вибору коефіцієнтів, повторного квантування та корекції помилок у DCT-, F5- і JPEG-орієнтованих схемах \cite{parsafar2024qft,maurya2023quantumdct,wang2023hyperchaotic,pan2022hiddenbit,sushil2024f5,chen2024f5robustness,duan2023jpeg,pan2024jpeg}. Піксельне порівняння після декодування не відтворює початкового механізму, тому його межею є втрата частини коефіцієнтних свідчень. Для власного рішення потрібне окреме частотне подання; роздільне оброблення просторового й JPEG-подань у \cite{he2025twobranch} підтверджує доцільність різної передобробки, але не визначає готової схеми їх інтеграції.

Механізм вейвлетного розміщення використовує багатомасштабний розклад: низькочастотна складова передає основну структуру, а деталізувальні піддіапазони зберігають локалізовані зміни різного масштабу та орієнтації. Генеративні, дифузійні й гібридні схеми навчають розподіл змін або формують приховане подання разом із носієм \cite{yao2024ganwavelet,yuan2022gansecurity,peng2023stegaddpm,ramandi2024vidagan,ma2023adversarial,ahmad2024cnndct}, тоді як стохастичне й багатошарове вбудовування змінює співвідношення місткості, непомітності та стійкості \cite{elhady2024dish,weng2024emd}. Спостережуваними залишаються локальні багатомасштабні залежності, але межею є перенесення на невідомий спосіб формування змін. Отже, власне рішення має зберігати деталізувальні піддіапазони й їхню просторову прив’язку, не приписуючи їм універсальної роздільної здатності.

Корисною властивістю вейвлетних коефіцієнтів є просторова локалізація деталізувальних компонентів. Wave-ViT використовує оборотне вейвлетне зменшення розмірності зі збереженням високочастотної інформації під час механізму уваги \cite{yao2022wavevit}. Ця архітектура не є стеганалізатором, тому її результат не слід тлумачити як підтвердження придатності до задачі виявлення. Вона лише обґрунтовує збереження деталізувальних піддіапазонів як загальний принцип побудови зорового представлення. Подібно, FcaNet інтерпретує глобальне усереднення як окремий частотний випадок і розширює його спектральними компонентами \cite{qin2021fcanet}. Для досліджуваної задачі це є підставою оцінювати канали за частотним вмістом, але не замінює емпіричної перевірки стеганалітичного модуля.

Спільною умовою просторового, частотного й вейвлетного розміщення є компроміс між місткістю, непомітністю та стійкістю прихованих даних. Збільшення кількості змінених елементів полегшує передавання більшого фрагмента, але за інших однакових умов посилює статистичне відхилення. Пристосування змін до текстури або країв зменшує зорову помітність, проте переносить інформативний слід у локальні залежності. Орієнтація на стійкість до повторного кодування змінює вибір компонентів і може зробити неінформативним детектор, налаштований на інший механізм. Огляди стеганографії та стеганалізу розглядають цю взаємозалежність як одну з причин постійної зміни ознак відомих способів \cite{muralidharan2022race,apau2024review,kheddar2024review}. Для вебресурсу до неї додається маршрут перетворень: однакове початкове вбудовування після масштабування, побудови ескізу або перекодування залишає інше доступне подання. Тому аналіз назви алгоритму чи одного показника якості не визначає очікуваної виявлюваності. Потрібно зіставляти область зміни, правило її вибору, формат носія та операції, виконані до моменту контролю. Такий підхід не передбачає універсальної ознаки, але дає змогу обґрунтовано формувати кілька взаємодоповнювальних джерел свідчень.

У відео механізм розміщення може змінювати окремі кадри, розподіляти дані між часовими позиціями або використовувати параметри стисненого потоку. Спостережуваними є локальні покадрові відхилення, їхня часова послідовність і кодекові параметри. Мале відхилення в кожному кадрі та коротка тривалість прояву обмежують випадкову вибірку й надмірне усереднення. Вимоги до одиниці висновку, часової прив’язки та правила відбору обґрунтовано під час порівняння відеометодів у підрозділі~\ref{sec:detection_methods_analysis}; тут суттєвим є лише розмежування трьох місць приховування.

Кодекове подання містить коефіцієнти дискретного косинусного та дискретного синусного перетворень (DCT/DST), параметри прогнозування, індекси кандидатів векторів руху й режими поділу блоків. Метод у \cite{zhang2022cnnhevc} орієнтований на сліди DCT/DST-вбудовування у стандарті високоефективного кодування відео (HEVC), підхід \cite{liu2021hevclocaloptimality} використовує порушення локальної оптимальності кандидатів руху, а в \cite{cao2026hevcpartition} аналізують режими поділу блоків перетворення. Ці джерела показують, що покадрове декодування вилучає частину початкових кодекових ознак; вони не доводять універсальності одного детектора для різних стандартів і способів вбудовування.

Структурне розміщення не потребує зміни пікселів або кадрів. Фізична байтова послідовність може містити область після логічного завершення JPEG чи PNG, допоміжний сегмент, приватне поле або вкладений фрагмент. Така область не є загрозою лише через своє розташування, оскільки окремі робочі процеси допускають розширення. Аномальність визначається невідповідністю політиці формату, ознаками іншого типу даних або зв’язком із несанкціонованою дією.

Тип допоміжної області можна оцінювати за внутрішніми байтовими закономірностями. FiFTy \cite{mittal2021fifty}, Byte2Vec \cite{haque2022byteembeddings} і ByteRCNN \cite{skracic2023bytercnn} спостерігають відповідно локальні, контекстні та послідовні властивості фрагмента, але сам визначений тип не встановлює його призначення. Дослідження витоку ескізів JPEG і карта трас походження зображення пов’язують похідні подання з історією перетворень \cite{benes2023thumbnails,kiltz2024trace}. Отже, структурне рішення має перевіряти не лише початкову сигнатуру, а й тип, позицію, допустимість та походження кожної знайденої області.

Багатоформатний об’єкт використовує допустимі пропуски, коментарі, зміщення або правила ігнорування невідомих даних, щоб залишатися коректним для кількох інтерпретаторів. Один компонент розпізнає зображення, а інший може знайти архівну, документну чи сценарну структуру. Дослідження \cite{koch2022polyglot} виявило обмеження засобів, що покладаються на одну сигнатуру або евристичний пошук вкладень. У \cite{koch2025polyglot} аналіз доповнено безпечною реконструкцією зображень, що переводить результат від простого сигналу тривоги до контрольованої дії над об’єктом.

Службові поля EXIF і XMP можуть містити текст, двійкові ескізи, ідентифікатори, дані про програмне забезпечення та розширення. Приховування в описових полях і вилучення EXIF з багатоформатних об’єктів розглянуто в \cite{fernando2024metadata,kalaimagal2024exif}, а документований ExifTool забезпечує відтворюваний розбір таких полів \cite{exiftool2026}. У штатному процесі службові дані підтримують сумісність, керування активами й цифрову криміналістику \cite{easttom2024windows}. Тому суцільне видалення службових полів може знищити потрібні відомості. Аналіз має встановлювати їхню синтаксичну коректність, семантичну відповідність і допустимість за політикою вебресурсу.

Небезпечний прояв може виникнути без відновлення текстової команди. Спеціально сформовані довжини, порядок блоків або параметри здатні взаємодіяти з помилкою декодера чи засобу перетворення. Відомості про уразливості ImageIO та ImageMagick показують, що така дія можлива на різних етапах інтерпретації мультимедійних даних \cite{opencve2023cve41064,apple2023cve41064,suse2023cve2157,li2023imagick}. Отже, аномальний код охоплює також небезпечні значення та структури, які спричиняють неочікувану послідовність операцій. Формальна коректність для одного засобу не гарантує безпечності для іншої реалізації.

Функціональна роль прихованої послідовності не завжди збігається зі способом її розміщення. У зоровому або службовому поданні може міститися безпосереднє програмне навантаження, однак та сама область може передавати коротку команду, ключ, адресу чи умову подальшого відновлення. Інший випадок становить спеціально сформоване значення, яке не потребує окремого вилучення, оскільки впливає на хід роботи вразливого обробника. Приклади приховування шкідливого коду, запускання за сигналом і доставляння даних через зображення підтверджують різні ролі носія в ланцюгу впливу \cite{kumagai2022malwareimage,almehmadi2022triggering,iglesias2023driveby}. Відомості про помилки мультимедійних засобів, своєю чергою, показують, що небезпечна семантика може бути закладена у структурний параметр, а не у відновлюваний текст \cite{opencve2023cve41064,suse2023cve2157}. Це розмежування не дає змоги встановити призначення фрагмента лише за його позицією. Воно визначає, які додаткові свідчення потрібні: вміст і тип області для навантаження, зв’язок із зовнішньою логікою для команди або вказівника, реакція конкретного засобу для небезпечного параметра. Отже, структурний розбір, аналіз спектральних і візуальних ознак та спостереження за обробленням відповідають на різні частини одного питання.

У середовищі вебсистеми прихована послідовність може виконувати три різні ролі. По-перше, вона може бути частиною доставляння програмного навантаження, яке зовнішня логіка зчитує з піксельного або службового подання \cite{iglesias2023driveby,asharani2024stegosploit,mohan2024stegomalware,badar2025stegomalware,huang2024shellcode,cassavia2022threats}. По-друге, носій може передавати умову або команду запускання, тому значення спостерігається лише після зіставлення з пов’язаним сценарієм \cite{almehmadi2022triggering,checkpoint2024v8,hp2022magniber,khan2023ransomware}. По-третє, спеціально сформований параметр може впливати на сам засіб оброблення без окремого відновлення текстової команди. Межа статичного аналізу полягає в тому, що він не підтверджує фактичну активацію жодної з цих ролей. Ізольоване виконання формує окреме джерело подієво-поведінкових свідчень \cite{abbadini2023sandboxing}, але не замінює встановлення області, типу й походження даних у носії. Тому власне рішення має зв’язувати статичне відхилення з подіями саме того сценарію, у якому оброблявся конкретний об’єкт.

Комбінований спосіб розподіляє взаємопов’язані частини між представленнями. Наприклад, пікселі можуть містити фрагмент послідовності, службове поле може зберігати ключ або вказівник, а додаткова байтова область доповнювати навантаження. Слабкість кожного окремого відхилення не усуває їхнього спільного значення. Саме тому потрібне мультимодальне представлення, яке зберігає належність ознак одному об’єкту та дає змогу оцінити їхню узгодженість.

Способи прихованого розміщення зіставлено в табл.~\ref{tab:hiding_methods} за рівнем, впливом штатних перетворень у вебсистемі і джерелами ознак. Характеристики є якісними: фактичне збереження даних залежить від формату, параметрів кодування та конкретного процесу оброблення.

\begin{table}[H]
\caption{Порівняльна характеристика способів прихованого розміщення аномальних кодів}
\label{tab:hiding_methods}
\centering
\scriptsize
\begin{tabularx}{\textwidth}{>{\raggedright\arraybackslash}p{0.17\textwidth}>{\raggedright\arraybackslash}p{0.23\textwidth}YY}
\toprule
\textbf{Рівень розміщення} & \textbf{Принцип приховування} & \textbf{Вплив штатного перетворення} & \textbf{Провідні джерела ознак} \\
\midrule
Просторовий & Контрольована зміна пікселів, каналів кольору або прозорості & Масштабування й кодування з втратами можуть руйнувати дані; безвтратне збереження переважно їх залишає & Бітові й піксельні статистики, локальні залишки, шумовий профіль \\
DCT-частотний & Зміна вибраних квантованих коефіцієнтів у блоках JPEG або відеокодека & Залежить від повторного квантування та якості кодування & Розподіли DCT, міжблокові залежності, частотні залишки \\
Вейвлетний & Розподілення даних між деталізувальними піддіапазонами різних масштабів & Частина ознак може зберігатися після помірних перетворень, але залежить від обраного розкладу & Вейвлетні коефіцієнти, багатомасштабні карти, локалізовані високочастотні зміни \\
Часовий і кодековий & Зміна кадрів, векторів руху, параметрів прогнозування чи поділу блоків & Перекодування може змінити кодекові ознаки; покадрове декодування втрачає частину інформації & Міжкадрові залежності, вектори руху, DCT/DST, режими HEVC \\
Структурний & Додаткові області, вкладені фрагменти, багатоформатна будова & Зберігається під час простого копіювання; може вилучатися контрольованою реконструкцією & Маркери, довжини, порядок блоків, типи байтових фрагментів, друга інтерпретація \\
Службовий & Розміщення даних у EXIF, XMP, коментарях та допоміжних полях & Очищення службових даних може вилучити навантаження, але не завжди допустиме політикою використання & Синтаксис, довжина, семантика й контекст споживання службових полів \\
Комбінований & Розподілення взаємопов’язаних частин між кількома рівнями & Окреме перетворення може зруйнувати один канал, не усунувши інші прояви & Узгоджена сукупність спектрально-зорових, структурних, часових і поведінкових ознак \\
\bottomrule
\end{tabularx}
\end{table}

Місце приховування визначає початковий набір потенційно інформативних ознак, але не однозначний висновок. Частотне вбудовування потребує аналізу коефіцієнтів і залишків, додаткова байтова область потребує структурного розбору, а дія під час оброблення потребує контролю подій. Значення кожної ознаки залежить від формату та політики: додаткові байти можуть бути допустимими, а високочастотний залишок може виникати через сенсор або повторне кодування.

Отже, зорові, частотні, вейвлетні, часові, кодекові, структурні та службові способи утворюють взаємопов’язану систему. Аналіз відомих методів має встановити, які представлення вони спостерігають, які перетворення враховують і чи здатні пов’язати властивості контейнера з реакцією вебсистеми.

\section{Аналіз відомих методів виявлення аномальних кодів у мультимедійних об’єктах}
\label{sec:detection_methods_analysis}

Відомі методи розв’язують різні задачі: установлюють наявність прихованих даних, оцінюють їхній обсяг, локалізують змінені області, визначають спосіб приховування або фіксують небезпечну дію під час оброблення. Ці результати не є взаємозамінними. Висока частка правильних двокласових рішень не підтверджує здатності до локалізації, а структурна невідповідність сама по собі не встановлює семантичного призначення знайденої області.

Для порівняння потрібно враховувати вид носія, рівень і місткість вбудовування, доступність початкового об’єкта, історію кодування, джерело даних, правило розподілу та тип рішення. Огляди засвідчують істотні відмінності між протоколами підготовки зображень \cite{muralidharan2022race,kheddar2024review}. Для відео додатково змінюються кодек, профіль, тривалість фрагмента й спосіб відбору кадрів \cite{bouzegza2022video}. Тому в цьому розділі зіставляються спостережувані рівні, механізми виділення ознак і встановлені обмеження, а числові показники мають оцінюватися за спільним протоколом.

Статистичні методи спираються на припущення, що вбудовування змінює природні розподіли значень, парність, локальні різниці, залежності сусідніх елементів або частоти коефіцієнтів. Гістограмні й бітові критерії є інтерпретованими та відносно маловитратними, проте їхня чутливість пов’язана з конкретним механізмом зміни. Адаптивне приховування в текстурних областях може залишати глобальний розподіл близьким до природного, тому відсутність вираженого відхилення не є підставою для автоматичного допуску.

Моделі локальних залежностей аналізують різниці, марковські переходи та співвиникнення квантованих залишків. Банки високочастотних фільтрів пригнічують семантичний зміст і підсилюють слабкі зміни, але фіксований набір фільтрів відображає лише закономірності, передбачені під час проєктування. Зміна формату, джерела чи способу вбудовування може знизити інформативність таких ознак. Межу між ручною моделлю залишків і навчуваним представленням розглянуто в оглядах та порівняльних дослідженнях \cite{kheddar2024review,muralidharan2022race,kuznetsov2024deep}.

Для JPEG статистики доцільно будувати безпосередньо за DCT-коефіцієнтами, їхніми гістограмами, знаками, парністю й міжблоковими залежностями. Це зберігає зв’язок із простором внесення даних до втрати інформації під час декодування. Водночас результат залежить від якості, повторного квантування, камери та редакторського процесу. Частотні статистики тому є окремим джерелом ознак, а не універсальним правилом для всіх носіїв.

Структурні методи перевіряють внутрішню граматику формату: межі сегментів, довжини, порядок блоків, контрольні суми, завершальні маркери, відповідність заявленого типу фактичній будові та семантику службових полів. Такий контроль виконується до повного декодування й може виявити області, які не впливають на зорове подання. Його обмеженням є підтримка допустимих розширень і невідомих блоків у самих форматах; надмірно жорстке правило здатне відхилити коректний рідкісний об’єкт.

Навчувана класифікація байтових фрагментів зменшує залежність від явних сигнатур, але спостерігає різні властивості. FiFTy використовує локальний вміст фрагмента без заголовка й відомостей про місце збереження \cite{mittal2021fifty}, Byte2Vec формує контекстне байтове подання \cite{haque2022byteembeddings}, а ByteRCNN поєднує локальні й послідовні залежності \cite{skracic2023bytercnn}. Спільною межею є відсутність семантики розташування: визначений тип не пояснює, чому фрагмент міститься у конкретній області мультимедійного контейнера. Тому у власному рішенні ці моделі можуть формувати структурні ознаки, але висновок потребує зіставлення типу з позицією, початковим записом і політикою формату.

Виявлення багатоформатних об’єктів безпосередньо спрямоване на неоднозначність інтерпретації. У \cite{koch2022polyglot} порівняно поширені засоби визначення типу й навчувані моделі, а в \cite{koch2025polyglot} аналіз доповнено зразками з ланцюгів атак і контрольованою реконструкцією. Переносимість такого детектора залежить від пар форматів і способів розташування фрагментів у навчальному наборі. Крім того, наявність другої структури не доводить її виконання в конкретній вебсистемі.

Структурний результат має різну доказову силу залежно від того, що саме встановлено. Порушення довжини або порядку елементів свідчить про невідповідність граматиці, але не визначає вміст нетипової області. Визначення типу вмісту уточнює ймовірний тип байтового фрагмента, проте не встановлює його призначення без урахування позиції та політики контейнера \cite{mittal2021fifty,haque2022byteembeddings,skracic2023bytercnn}. Виявлення другої форматної інтерпретації посилює підставу для поглибленої перевірки, однак коректні розширення й допоміжні дані можуть мати законне походження. Навіть контрольована реконструкція показує можливість вилучити неоднозначну частину, а не доводить, що вона була використана у конкретній вебсистемі \cite{koch2022polyglot,koch2025polyglot}. Отже, структурна гілка повинна зберігати вид порушення, байтове розташування й результат визначення типу як окремі свідчення. Їхнє значення уточнюється спектрально-зоровими властивостями та подіями оброблення, а не замінюється загальною позначкою некоректності.

Спектрально-зорові методи спрямовано на слабкі зміни зорового подання. Огляди 2023--2025 років групують їх за переходом від ручних статистик до глибоких, генеративних і предметно орієнтованих моделей та одночасно вказують на залежність результату від носія, алгоритму й протоколу \cite{rustad2023survey,wani2023review,luo2024survey,michaylov2025forensic}. Спостережуваним для цієї групи є низькоамплітудний локальний слід, тому типовий механізм містить високочастотну або залишкову передобробку, виділення ознак і класифікацію. Межа полягає в тому, що представлення, корисне для розпізнавання предметного змісту, може пригнічувати такий слід як шум. Отже, власний метод має зберігати детерміновану залишкову основу незалежно від наявності повного навчуваного кодувальника.

За механізмом формування ознак архітектури утворюють три підгрупи. Локально-залишкові моделі використовують згорткове агрегування, SRM-передобробку й залишкові блоки \cite{fu2022ccnet,bravoortiz2024cvtstegonet,ntivuguruzwa2023cnn}; вони спостерігають слабкі локальні зміни, але залежать від закладеного масштабу фільтрів. Моделі уваги й текстурного контексту доповнюють локальні залишки глобальними залежностями \cite{xie2023selfattention,zhou2024texturetransformer}, однак їхня перевага залежить від навчального корпусу. Предметно орієнтовані та ресурсно скорочені рішення використовують високовимірне JPEG-подання, великі еталонні масиви, скорочення параметрів, еволюційний добір або злиття рішень \cite{xu2025jpegmap,gao2025recognition,huang2024sasrnet,ma2025evolutionary,swarnkar2024fusion}. Спільною межею всіх трьох підгруп є орієнтація на зоровий носій без повного розбору контейнера. Тому у власному рішенні навчуване зорове подання може лише доповнювати структурну гілку, а не підміняти її.

Механізм відбору областей і носіїв визначає, які відхилення взагалі потраплять до аналізу. Високоентропійний відбір зменшує обсяг обчислень і зосереджується на складних фрагментах \cite{agarwal2024entropy}, але може пропустити гладку або суто структурну область. Адаптивна фільтрація залишків із контролем добору початкових зображень \cite{ma2023coverselection}, вибір домінантних ознак із компенсацією \cite{yu2024dominantfeatures}, огляд вибору носіїв і графове групування \cite{azame2025coverselection,dehdar2023graph} показують іншу межу: модель може засвоїти властивості джерела замість сліду приховування. Отже, власне рішення має зберігати правило відбору й походження носія, а пропущену область не тлумачити як підтвердження норми.

Для експлуатації важливі також обчислювальні витрати. LWENet використовує полегшені залишкові блоки, роздільні згортки та кілька способів глобального агрегування \cite{weng2022lwenet}. Блокове проріджування в \cite{hong2023blockpruning} зберігає значущі ранні низькорівневі блоки. Такі результати обґрунтовують каскадне спрямування частини об’єктів на поглиблений аналіз, але скорочення потрібно перевіряти окремо для кожної модальності.

За типом виходу спектрально-зорові методи формують різні, не взаємозамінні результати. Кількісний стеганаліз оцінює обсяг прихованих даних \cite{singhal2024ustegdse}, локалізація визначає змінені пікселі \cite{ntivuguruzwa2023localization}, а класифікація з реконструкцією намагається відновити приховане зображення \cite{lai2022ffrnet}. SPAM-ознаки, несумісні JPEG-блоки й неузгодженість класифікаторів дають локальні пояснювані свідчення \cite{hemalatha2023spam,levecque2024jpeg,lerch2023actor}; Aletheia забезпечує відкриту основу для частини аналізаторів \cite{lerch2024aletheia}. Межею є залежність виходу від механізму приховування: неможливість реконструкції не заперечує виявленого відхилення, а двокласова оцінка не забезпечує локалізації. Тому власне рішення має зберігати тип виходу й координати свідчення окремо від підсумкового класу.

Відеометоди можна розрізнити за покадровим, часовим і кодеково-структурним поданням. Покадровий аналіз повторно використовує модель для зображень, але не охоплює розподілене вбудовування та параметри стисненого потоку. Часове агрегування зберігає послідовність, однак потребує такої довжини фрагмента й правила об’єднання, які не згладжують короткі відхилення. Кодеково-структурні ознаки є інформативними для відповідного механізму, але залежать від стандарту й реалізації кодера.

Локальну оптимальність векторів руху використано в \cite{liu2021hevclocaloptimality}, а узагальнення цього принципу подано в \cite{zhai2025generalized}. Метод централізованої помилки й уваги аналізує стеганографію коефіцієнтів перетворення з урахуванням прогнозування HEVC \cite{dai2025centralized}; режими поділу блоків перетворення досліджено в \cite{cao2026hevcpartition}. Чутливість до векторів руху не гарантує чутливості до коефіцієнтів, а детектор кадрових залишків може не спостерігати рішення кодера. Це обмежує пряме перенесення результатів між відеосценаріями.

Одиницею висновку для відео має залишатися цілісний об’єкт, хоча окремі свідчення виникають на рівні кадру, часового інтервалу або елемента стисненого потоку. Якщо зміна зосереджена в короткому фрагменті, випадковий відбір кадрів може не охопити її, а усереднення послабити локальне відхилення. Аналіз кожного кадру, навпаки, збільшує витрати й створює залежні спостереження, які не можна трактувати як окремі об’єкти. Кодекові ознаки доповнюють кадрове подання, але їхня відсутність після декодування означає втрату джерела даних, а не нормальний стан. Огляд відеостеганалізу та спеціалізовані методи для векторів руху, коефіцієнтів і режимів поділу підтверджують відмінність цих механізмів \cite{bouzegza2022video,liu2021hevclocaloptimality,dai2025centralized,cao2026hevcpartition}. Тому коректний аналіз має зберігати часову позицію кожної ознаки, правило відбору та зв’язок із контейнером. Така простежуваність дає змогу відокремити відсутність спостереження від відсутності прихованої зміни й визначає межі подальшого агрегування.

Поведінкові методи спостерігають процес контрольованого відкриття або перетворення. Журнал може містити операції читання й запису, мережеві звернення, створення процесів, події пам’яті та виклики прикладних засобів. Нормальний профіль залежить від сценарію: створення мініатюри й перекодування відео закономірно породжують різні події та витрати. Окрема подія тому не є універсальним доказом; потрібні послідовність, прив’язка до об’єкта й зіставлення з контрольним запуском.

Робота \cite{iglesias2023driveby} обґрунтовує спільний розгляд прихованих даних і можливої дії у середовищі вебсистеми. Однак відсутність активації в обраному сценарії може спричинити пропуск, а фонова діяльність засобів перетворення може створити хибну тривогу. Поведінкові події мають уточнювати структурні та спектрально-зорові ознаки того самого об’єкта, а не утворювати незалежний безумовний бал.

Причинна прив’язка поведінкової події потребує однакових умов для досліджуваного й контрольного оброблення. Версія засобу, початковий стан середовища, права, мережеві обмеження та фонові служби впливають на журнал незалежно від властивостей носія. Порівняння має стосуватися не лише наявності події, а й її місця в послідовності, часу відносно відкриття об’єкта та повторюваності за незмінного сценарію. Поведінкові шаблони допомагають відокремити закономірну роботу компонента від нетипового ланцюга \cite{denysiuk2024behavior}, а ізольоване середовище обмежує сторонній вплив під час спостереження \cite{abbadini2023sandboxing}. Проте контрольоване оброблення не гарантує активації прихованої логіки: умова активації може залежати від дії користувача, зовнішнього ресурсу або іншого стану, якого немає у вибраному сценарії. Через це відсутність події слід реєструвати разом з умовами оброблення, а не перетворювати на безумовне свідчення норми. Підстава для інтерпретації посилюється, коли нетипова послідовність має часовий зв’язок з обробленням конкретного носія та узгоджується з його структурним або спектрально-зоровим відхиленням.

Групування методів за джерелами даних і характером рішення наведено на рис.~\ref{fig:detection_method_classification}. Схема показує взаємодоповнення груп, а не ієрархію їхньої взаємозамінності.

\begin{figure}[htbp]
\centering
\begin{tikzpicture}[x=1cm,y=1cm,classification group/.style={dissertation block,text width=2.60cm,font=\footnotesize\setstretch{1}}]
\node[dissertation block,text width=3.50cm] (root) at (0,3) {Методи виявлення};
\node[classification group] (stat) at (-6.5,0) {Статистичні};
\node[classification group] (str) at (-3.9,0) {Структурні};
\node[classification group] (sv) at (-1.3,0) {Піксельні та\\спектральні};
\node[classification group] (video) at (1.3,0) {Ознаки\\закодованого\\відеопотоку};
\node[classification group] (beh) at (3.9,0) {Поведінкові};
\node[classification group] (multi) at (6.5,0) {Узгодження\\різнорідних\\даних};
\node[dissertation block,text width=9.20cm] (gap) at (0,-3) {Наукова прогалина: відсутність узгодженого подання\\рівнів мультимедійного об’єкта та реакції вебсистеми};
\draw[thick] (root.south) -- (0,1.55);
\draw[thick] (-6.5,1.55) -- (6.5,1.55);
\draw[dissertation arrow] (-6.5,1.55) -- (stat.north);
\draw[dissertation arrow] (-3.9,1.55) -- (str.north);
\draw[dissertation arrow] (-1.3,1.55) -- (sv.north);
\draw[dissertation arrow] (1.3,1.55) -- (video.north);
\draw[dissertation arrow] (3.9,1.55) -- (beh.north);
\draw[dissertation arrow] (6.5,1.55) -- (multi.north);
\draw[thick] (stat.south) -- (-6.5,-1.55);
\draw[thick] (str.south) -- (-3.9,-1.55);
\draw[thick] (sv.south) -- (-1.3,-1.55);
\draw[thick] (video.south) -- (1.3,-1.55);
\draw[thick] (beh.south) -- (3.9,-1.55);
\draw[thick] (multi.south) -- (6.5,-1.55);
\draw[thick] (-6.5,-1.55) -- (6.5,-1.55);
\draw[dissertation arrow] (0,-1.55) -- (gap.north);
\end{tikzpicture}
\caption{Групи методів виявлення ознак аномального вмісту за джерелами даних}
\label{fig:detection_method_classification}
\end{figure}


Статистичні та структурні методи переважно формують пояснювані ознаки до декодування або відразу після нього. Спектрально-зорові мережі навчають подання слабких змін, а поведінковий аналіз додає дані, яких немає в контейнері. Узгодження має зберігати сильне підтвердження окремого джерела та окремо відображати конфлікт між модальностями, а не приховувати його усередненням.

Порівняльну характеристику наведено в табл.~\ref{tab:detection_comparison}. Провідне обмеження означає властивість, яку потрібно компенсувати іншими даними, а не абсолютну нездатність групи методів.

\begin{table}[H]
\caption{Порівняння груп методів виявлення аномальних кодів у мультимедійних об’єктах}
\label{tab:detection_comparison}
\centering
\scriptsize
\begin{tabularx}{\textwidth}{>{\raggedright\arraybackslash}p{0.15\textwidth}>{\raggedright\arraybackslash}p{0.20\textwidth}>{\raggedright\arraybackslash}p{0.21\textwidth}YY}
\toprule
\textbf{Група методів} & \textbf{Вхідні дані} & \textbf{Сильна сторона} & \textbf{Провідне обмеження} & \textbf{Роль в аналізі} \\
\midrule
Статистичні & Піксельні й частотні розподіли, залишки, співвиникнення & Інтерпретовані низькорівневі індикатори й малі часові витрати & Залежність від закладеної моделі носія та способу вбудовування & Раннє оцінювання й допоміжні ознаки для зорової гілки \\
Структурні & Повна байтова будова, маркери, блоки, службові поля & Виявлення додаткових областей і форматної неоднозначності до повного декодування & Не спостерігають слабких змін, що не порушують будови контейнера & Контроль формальної коректності та формування структурного представлення \\
Спектрально-зорові & Пікселі, DCT, вейвлети, залишкові карти, кадри & Чутливість до слабких і просторово розподілених змін & Зміщення джерела носіїв і відсутність даних про службові області & Аналіз прихованих змін у зображенні й кадрах \\
Відеокодекові & Вектори руху, режими прогнозування й поділу, DCT/DST & Спостереження слідів, які зникають після покадрового декодування & Прив’язка до кодека, профілю та цільового способу вбудовування & Аналіз стисненого відеопотоку \\
Поведінкові & Події контрольованого декодування, перетворення й відображення & Фіксують можливий прояв прихованого впливу під час оброблення & Залежність від умов активації, середовища та фонового профілю & Перевірка зв’язку між властивостями об’єкта і реакцією вебсистеми \\
Мультимодальні & Узгоджені ознаки кількох рівнів & Підтвердження розподілених проявів і пояснення конфліктів між джерелами & Потреба в синхронізації, калібруванні та даних із повними модальностями & Формування інтегрованої оцінки ризику й підтримка рішення \\
\bottomrule
\end{tabularx}
\end{table}

Переносимість обмежує зміщення джерела даних. Контрольовані набори часто містять вихідні й змінені зображення однакового походження, тоді як потік вебресурсу поєднує камери, редактори, генератори та повторне кодування. Вибір носіїв у \cite{ma2023coverselection} і випробування на ALASKA\#2 у \cite{xie2023selfattention} частково враховують цю проблему, а набори природних фото- й відеозаписів із багатьох пристроїв, як-от VISION~\cite{shullani2017vision}, дають матеріал для оцінювання переносимості на реальні джерела; проте жоден із них не замінює перевірки на даних цільового середовища. Якщо класи систематично відрізняються джерелом, модель може розпізнавати джерело замість сліду вбудовування.

Зміщення може виникати на кожному етапі формування носія. Пристрій або генератор визначає початковий шумовий профіль; редакторська обробка змінює локальні залежності; формат і параметри кодування впливають на коефіцієнти; спосіб приховування задає розташування та інтенсивність сліду. Окремим чинником є склад потоку, оскільки частка форматів і джерел у вебресурсі може відрізнятися від навчальної вибірки. Огляди вибору носіїв показують, що джерельні властивості потрібно контролювати так само, як механізм вбудовування \cite{azame2025coverselection,ma2023coverselection}. Інакше супутня відмінність стає простішою основою рішення, ніж слабка ознака прихованої послідовності. Подібний ризик створюють споріднені похідні копії: мініатюра й початковий об’єкт мають спільний зміст та історію, тому їх розподілення між різними частинами набору не забезпечує незалежної перевірки \cite{kiltz2024trace}. Аналіз переносимості через це має розрізняти зміну джерела, зміну перетворення та зміну способу приховування. Зниження класифікаційного показника в такому сценарії не дає змоги без додаткового розбору приписати помилку одному чиннику, але чітке групування умов показує межу, у якій ознаки залишаються придатними.

Неповнота модальностей є штатною умовою. EXIF або XMP можуть бути відсутні, параметри кодека можуть бути недоступні, а контрольований запуск може не породити подій. Заповнення пропусків нулями здатне перетворити сам факт відсутності на небажану ознаку класу. Інформаційна технологія виявлення аномальних кодів у фото- та відеооб’єктах вебресурсів повинна зберігати маску доступності, надійність джерел і можливість рішення за неповним набором ознак.

Числові результати потрібно співвідносити зі складністю протоколу. Закрите випробування з відомими алгоритмом, місткістю та джерелом перевіряє відтворення навченої межі, а невідомий алгоритм або нове джерело перевіряють переносимість. Масштабування, повторне кодування й очищення службових полів утворюють додаткові сценарії для вебресурсів. Показники цих сценаріїв відповідають різним дослідницьким питанням і не можуть тлумачитися як результати однієї задачі.

Одиницею розподілу має бути початковий носій або повний відеофрагмент разом із похідними поданнями. Якщо споріднені копії або кадри одного відео потрапляють до навчальної й тестової частин, модель може розпізнавати зміст, пристрій або часову близькість замість ознаки приховування.

Вихідна оцінка мережі не дорівнює емпіричній частоті загрози. Зміна джерела, формату або частки класів зміщує робочу точку, а поріг, що максимізує частку правильних рішень на збалансованій вибірці, може створити неприйнятну кількість тривог у робочому потоці вебресурсу. Необхідно зберігати неперервну оцінку, умови її отримання та надійність, а робочу точку пов’язувати з вартістю пропуску й пропускною здатністю поглибленої перевірки.

Пояснення потребує узгодження координат. Карта уваги характеризує чутливість моделі, але не доводить розташування коду; байтовий інтервал точно локалізує структурне відхилення, але не завжди визначає його призначення; подія показує реакцію, але не доводить причинного зв’язку. Результат має містити тип ознаки, просторову, часову або байтову локалізацію, надійність джерела й підтвердження іншими даними.

Ресурсну придатність визначає повний критичний шлях, а не кількість параметрів окремої мережі. Потрібно враховувати структурний розбір, декодування, відбір кадрів, передавання даних і контрольований запуск; вимірювати медіану й верхні квантили часу, пікову пам’ять, пропускну здатність і частку поглиблених перевірок. Полегшені архітектури \cite{weng2022lwenet,hong2023blockpruning} слід оцінювати за впливом на повне рішення.

Рівень інтеграції також змінює зміст результату. Статистичний залишок можна подавати як вхідний канал, просторові та JPEG-ознаки можна обробляти паралельно \cite{he2025twobranch}, а структурне й поведінкове рішення можна об’єднувати наприкінці. Раннє злиття потребує сумісних розмірностей, пізнє може втратити зв’язок між конкретною областю та подією. Поведінкове узгодження ознак різного походження має враховувати маски доступності, надійність і походження ознак.

Критерії інтерпретації сучасних методів 2021--2026~рр. наведено в табл.~\ref{tab:evidence_quality}. Вони відокремлюють властивість детектора від умов, за яких її спостерігали. Для аналізу робочої точки розрізняють криву робочих характеристик приймача (ROC) та криву «точність--повнота» (PR).

\begin{table}[H]
\caption{Критерії доказовості результатів методів виявлення аномальних кодів}
\label{tab:evidence_quality}
\centering
\scriptsize
\begin{tabularx}{\textwidth}{>{\raggedright\arraybackslash}p{0.20\textwidth}>{\raggedright\arraybackslash}p{0.27\textwidth}YY}
\toprule
\textbf{Критерій} & \textbf{Що потрібно контролювати} & \textbf{Наслідок порушення} & \textbf{Вимога до дисертаційного протоколу} \\
\midrule
Незалежність даних & Групування початкового об’єкта з усіма його цілеспрямовано зміненими й похідними копіями; групування кадрів одного відео & Витік спорідненого змісту та завищення показників & Розподіл за початковими носіями до формування похідних прикладів \\
Переносимість & Невідомі алгоритми, джерела, формати, кодеки та параметри перетворення & Підміна виявлення загрози розпізнаванням умов набору & Окремі сценарії внутрішнього й зовнішнього випробування \\
Калібрування & Відповідність оцінки емпіричному ризику, стабільність порога за зміни частки класів & Надмірна кількість хибних тривог або пропусків у робочому потоці & Аналіз ROC- і PR-кривих, матриця помилок і обґрунтування робочої точки \\
Локалізація та пояснення & Координати, байтові інтервали, джерела ознак і причинний зв’язок із подіями & Висновок неможливо перевірити або використати для реагування & Збереження внесків модальностей і проміжних результатів аналізу \\
Ресурсна придатність & Повний час послідовності оброблення, пікова пам’ять, пропускна здатність, частка поглиблених перевірок & Швидкість окремої моделі не відображає затримки технології & Вимірювання етапів і загального критичного шляху \\
Відтворюваність & Версії, параметри, випадкові початкові значення, правила підготовки та метрики & Неможливість незалежно повторити або пояснити відмінний результат & Журнал конфігурації та однозначний опис усіх процедур \\
\bottomrule
\end{tabularx}
\end{table}

Високий показник на спеціалізованому наборі підтверджує працездатність лише в межах відповідного протоколу. Результат структурного аналізу багатоформатних об’єктів не переноситься на слабке адаптивне вбудовування, а показник HEVC-детектора не характеризує контейнерну аномалію MP4. Тому порівнювати потрібно не абстрактну перевагу класифікаторів, а здатність сформувати, узгодити й пояснити взаємодоповнювальні свідчення за визначених умов застосування.

У межах проаналізованого корпусу джерел наукову прогалину визначено як недостатню опрацьованість узгодженого представлення, яке пов’язує зорові, частотні, вейвлетні, структурні, службові, часові та подієво-поведінкові ознаки з маршрутом конкретного об’єкта у вебсистемі. Відомі мережі виявляють слабкі просторові зміни, відеометоди аналізують вектори руху й коефіцієнти HEVC, а структурні засоби знаходять нетипові фрагменти та багатоформатні об’єкти. Проте локальні оцінки переважно залишаються непорівнюваними або об’єднуються без збереження причинного зв’язку.

Цю прогалину визначають взаємні обмеження груп. Спектрально-зоровий аналіз не спостерігає даних після завершального маркера та другу форматну інтерпретацію. Структурний контроль не визначає слабке адаптивне вбудовування в природних областях зображення й не показує реакцію декодера. Поведінкова подія без прив’язки до носія не встановлює джерельну область і залежить від фонового процесу. Отже, потрібні метод спектрально-зорового та структурного аналізу фото- і відеооб’єктів вебресурсів та метод поведінкового узгодження ознак різного походження під час виявлення аномальних кодів у фото- та відеооб’єктах вебресурсів. Подальша постановка має визначити їхні входи, виходи, межі й місце в інформаційній технології виявлення аномальних кодів у фото- та відеооб’єктах вебресурсів.

Дослідницьке припущення, що випливає з установленої прогалини, полягає в такому: якщо статичні та подієво-поведінкові свідчення пов’язувати за спільним об’єктом, походженням, локалізацією й доступністю, то підсумкове представлення зберігатиме підстави рішення за неповного набору спостережень. Це припущення не означає наперед заданої переваги повної навчуваної конфігурації. Його мінімальним підтвердженням мають бути відтворюваний зв’язок між входами й виходом, окреме оцінювання доступних груп ознак та відсутність підміни пропущеної модальності нульовим свідченням. Класифікаційну перевагу можна встановлювати лише за незалежного розподілу й повних пооб’єктних результатів; таких даних у наявному корпусі немає.

\section{Постановка задачі дослідження та обґрунтування напряму формування інформаційної технології}
\label{sec:research_problem}

Задача виявлення полягає не в пошуку однієї наперед заданої сигнатури, а в оцінюванні узгодженості свідчень можливого прихованого впливу за неповними й різнорідними спостереженнями. Вхід становлять фотооб’єкт або відеооб’єкт, його службові дані та, за наявності, журнал контрольованого оброблення. Носій може бути формально коректним і зорово правдоподібним, а зміни можуть міститися у піксельному, частотному, вейвлетному, часовому, структурному чи службовому представленні. Вихід має містити пороговий висновок щодо достатності ознак аномального коду, можливий тип прихованої загрози, оцінку ризику, локалізаційні дані та пояснення за джерелами підтвердження.

Мультимедійний об’єкт розглядається як носій прихованого впливу, а вебсистема як середовище його приймання, перевірки, декодування, перетворення, зберігання та відображення. Дослідження не має на меті універсальну класифікацію всіх видів шкідливого програмного забезпечення та не замінює міжмережеві екрани, антивірусні шлюзи чи засоби ізольованого виконання. Його межа визначається спеціалізованим аналізом фото- та відеооб’єктів із прив’язкою властивостей носія до проявів оброблення.

Із цієї межі випливають вимоги до одиниці аналізу та простежуваності. Фотооб’єкт або повний відеофрагмент має отримувати незмінний ідентифікатор до декодування, а всі похідні кадри, мініатюри, структурні записи й журнали повинні зберігати зв’язок із цим носієм. Початкове байтове подання і безпечно декодоване подання виконують різні функції: перше зберігає контейнерні свідчення, друге відкриває зорові й спектральні залежності. Заміна одного подання іншим створює неповноту, яку неможливо усунути лише складнішим класифікатором. Тому результат має вказувати походження кожної ознаки та перетворення, після якого її отримано.

Доступність службових, кодекових і подієво-поведінкових даних залежить від формату та сценарію. Їхня відсутність повинна залишатися окремим станом спостереження, а не неявно прирівнюватися до нульового відхилення. За неповного набору свідчень допустимий висновок має відображати цю межу й за потреби спрямовувати носій на додаткову перевірку. Водночас оцінка властивостей об’єкта не повинна автоматично визначати однакову дію для всіх вебсистем. Політика допуску залежить від призначення ресурсу, можливості безпечної реконструкції та вартості помилкового блокування. Отже, предметом дослідження є формування обґрунтованого представлення загрози й підтримка рішення, тоді як конкретна адміністративна дія залишається результатом застосування цього представлення в заданій політиці.

Вихід аналізу також потребує розмежування за доказовим змістом. Пороговий висновок відповідає на питання про достатність сукупності свідчень, припущення щодо типу загрози пояснює можливий механізм, оцінка ризику враховує силу та узгодженість ознак, а локалізація вказує на область, яка потребує перевірки. Жоден із цих складників не повинен автоматично підміняти інші. Висока оцінка спектрального відхилення без структурної або подієво-поведінкової підтримки не встановлює функціональної ролі прихованих даних. Точне положення додаткової байтової області не доводить її активації. Подія без часової та джерельної прив’язки не дає надійної локалізації в носії. Тому пояснення має зберігати окремі підстави, їхню доступність і взаємне підтвердження. Така форма результату дає змогу перевірити рішення та вибрати наступну дію без приписування моделі можливостей, яких фактично не оцінювали.

Попередні публікації автора розглянуто лише як тематичні передумови, а не як незалежне підтвердження результатів цієї дисертації. Праці щодо виявлення стеганографічних змін і команд, переданих через зображення, безпосередньо стосуються формування статичних ознак і постановки задачі \cite{denysiuk2023steganalysis,denysiuk2025commands}. Дослідження поведінкових шаблонів і паралельних процесів пов’язані з організацією спостереження за послідовністю подій \cite{denysiuk2024behavior,denysiuk2025parallel}. Публікації про програмні імпланти, ботмережі, приманки та вторгнення описують ширший контекст контролю реакції середовища \cite{denysiuk2024implantscorporate,denysiuk2024botnets,denysiuk2024decoysystem,denysiuk2024systemdecoys,denysiuk2024intrusions}, але самі по собі не підтверджують мультимедійний метод. Ці праці окреслюють тематичний контекст дослідження й не використовуються в дисертації як взаємозамінні детектори; відомості про внесок здобувача у співавторські результати наведено у списку праць здобувача та відомостях про особистий внесок.

Збереження незмінного початкового мультимедійного об’єкта має передувати операціям, здатним змінити структуру. Паралельно потрібно формувати безпечно декодоване зорове подання. Така організація дає методу спектрально-зорового та структурного аналізу фото- і відеооб’єктів вебресурсів змогу опрацьовувати зв’язані представлення того самого носія без підміни структурної перевірки декодуванням.

Для відео необхідно зберігати часові мітки, параметри контейнера й доступні кодекові ознаки. Репрезентативний відбір кадрів має обмежувати обчислювальні витрати, але не згладжувати короткі відхилення. Покадрові результати повинні залишатися пов’язаними з часовим інтервалом і структурою вихідного відеоконтейнера.

Поведінковий прояв має контекстний зміст. Мережева операція, тимчасовий об’єкт або підвищене використання пам’яті можуть бути штатними для одного сценарію та нетиповими для іншого. Метод поведінкового узгодження ознак різного походження під час виявлення аномальних кодів у фото- та відеооб’єктах вебресурсів повинен зіставляти послідовність подій зі структурними й спектрально-зоровими відхиленнями та уточнювати ступінь підтримки гіпотези про загрозу.

Відсутність модальності не можна тлумачити як норму або аномалію. Потрібно зберігати маски доступності й оцінки надійності для службових, кодекових і подієво-поведінкових даних, а за недостатності свідчень підтримувати додаткову перевірку. Це відрізняє роботу з потоком вебресурсу від лабораторної класифікації з повним вектором для кожного запису.

Реагування має залежати від узгодженого ризику та підстав рішення. Структурна невідповідність може вести до додаткової перевірки чи безпечної реконструкції, а зовнішня коректність формату не гарантує допуску за наявності взаємопов’язаних слабких ознак. Потрібно підтримувати допуск, додаткову перевірку, реконструкцію або блокування та реєструвати підстави вибраної дії.

Предметну основу задає модель мультимодального представлення прихованої загрози. Метод спектрально-зорового та структурного аналізу фото- і відеооб’єктів вебресурсів формує передбачений моделлю статичний опис, а метод поведінкового узгодження ознак різного походження пов’язує його з реакцією середовища й формує інтегроване представлення. Інформаційна технологія виявлення аномальних кодів у фото- та відеооб’єктах вебресурсів організовує їх використання у п’ять етапів: попереднє оброблення, виділення ознак, інтеграція ознак, прийняття рішення та реєстрація результатів.

Зв’язок установлених обмежень із моделлю мультимодального представлення прихованої загрози, методом спектрально-зорового та структурного аналізу фото- і відеооб’єктів вебресурсів, методом поведінкового узгодження ознак різного походження та інформаційною технологією виявлення аномальних кодів у фото- та відеооб’єктах вебресурсів подано на рис.~\ref{fig:research_logic}.

\begin{figure}[htbp]
\centering
\begin{tikzpicture}[x=1cm,y=1cm,dissertation block/.append style={font=\small\setstretch{1}}]
\node[dissertation block,text width=4.20cm] (lim) at (0,5.5) {Обмеження відомих підходів};
\node[dissertation block,text width=9.00cm] (model) at (0,3.5) {Модель мультимодального представлення\\прихованої загрози};
\node[dissertation block,text width=5.90cm] (svs) at (-3.5,0.4) {Метод аналізу спектральних,\\піксельних і структурних ознак};
\node[dissertation block,text width=5.90cm] (beh) at (3.5,0.4) {Метод поведінкового\\узгодження ознак\\різного походження};
\node[dissertation block,text width=9.00cm] (tech) at (0,-2.8) {Інформаційна технологія виявлення ознак аномального вмісту\\в мультимедійних даних вебсередовища};
\node[dissertation block,text width=5.20cm] (exp) at (0,-5.1) {Експериментальне оцінювання\\та визначення доказових меж};
\draw[dissertation arrow] (lim) -- (model);
\draw[dissertation arrow] (model) -- (svs);
\draw[dissertation arrow] (model) -- (beh);
\draw[dissertation arrow] (svs) -- (tech);
\draw[dissertation arrow] (beh) -- (tech);
\draw[dissertation arrow] (tech) -- (exp);
\end{tikzpicture}
\caption{Логіка переходу від установлених обмежень до експериментального оцінювання}
\label{fig:research_logic}
\end{figure}


Відповідність установлених прогалин і необхідних складників дослідження наведено в табл.~\ref{tab:gap_results}.

\begin{table}[H]
\caption{Відповідність установлених прогалин і необхідних складників дослідження}
\label{tab:gap_results}
\centering
\scriptsize
\begin{tabularx}{\textwidth}{>{\raggedright\arraybackslash}p{0.25\textwidth}Y>{\raggedright\arraybackslash}p{0.13\textwidth}}
\toprule
Встановлена прогалина & Необхідний складник & Місце розкриття \\
\midrule
Роз’єднане подання ознак різних рівнів & Модель мультимодального представлення прихованої загрози & Розділ 2 \\
Ізольований зоровий або контейнерний аналіз & Метод спектрально-зорового та структурного аналізу фото- і відеооб’єктів вебресурсів & Розділ 3 \\
Відсутність зв’язку між носієм і реакцією середовища & Метод поведінкового узгодження ознак різного походження під час виявлення аномальних кодів у фото- та відеооб’єктах вебресурсів & Розділ 3 \\
Фрагментарне використання засобів контролю & Інформаційна технологія виявлення аномальних кодів у фото- та відеооб’єктах вебресурсів & Розділ 3 \\
Непорівнювані експериментальні твердження & Протокол оцінювання з компонентним вилученням, метриками помилок і часу та явними доказовими обмеженнями & Розділ 4 \\
\bottomrule
\end{tabularx}
\end{table}

Перехід від прогалин до результатів потребує також указати мінімальне свідчення й фактичний статус кожного складника:

\begin{enumerate}[label=\arabic*)]
\item для моделі мультимодального представлення прихованої загрози мінімальним свідченням є однозначний контракт входів і виходу зі спільним ідентифікатором, маскою доступності, локалізаціями та походженням ознак; модель формалізовано в розділі~\ref{chap:models_formalization}, а межі її експериментальної опори, зокрема неперевіреність повного мультимодального контуру на незалежних даних, установлено в розділі~\ref{chap:experiments};

\item для методу спектрально-зорового та структурного аналізу мінімальним свідченням є роздільне відтворення двох статичних гілок, їхніх локалізацій і внесків у спільне рішення; контрольний опис прогону, проаналізований у розділі~\ref{chap:experiments}, містить агреговані результати прототипної конфігурації для PNG і MP4, але не містить незалежного розподілу, пооб’єктних прогнозів і повної перевірки JPEG та WebP;

\item для методу поведінкового узгодження ознак різного походження мінімальним свідченням є причинно пов’язаний журнал того самого об’єкта та відтворюване зіставлення подій зі статичними локалізаціями; правилову конфігурацію формалізовано в розділі~\ref{chap:methods_technology}, тоді як повних журналів, навченого кодувальника послідовності й перевіреної перехресної уваги в наявному корпусі немає;

\item для інформаційної технології мінімальним свідченням є наскрізний запис п’яти етапів із незмінним входом, аналітичним станом, дією та аудитним записом; у наявних матеріалах зафіксовано час чотирьох операцій, а не повного п’ятиетапного контуру, що показано в розділі~\ref{chap:experiments};

\item для експериментального протоколу мінімальним свідченням є групово незалежний розподіл, пооб’єктні оцінки, матриця помилок, дані для аналізу робочих характеристик приймача та повний журнал конфігурації; доступна матриця дає змогу арифметично перевірити частину показників, але не забезпечує незалежного порівняння або повторного обчислення всіх заявлених характеристик.
\end{enumerate}

Перевірюваність обраного напряму потребує розмежування функцій його складників. Модель мультимодального представлення прихованої загрози має описувати склад, походження, доступність і взаємний зв’язок ознак, але сама по собі не підміняти процедуру їх вилучення. Метод спектрально-зорового та структурного аналізу фото- і відеооб’єктів вебресурсів повинен формувати простежуваний статичний опис із незмінного контейнера та декодованого подання. Метод поведінкового узгодження ознак різного походження має встановлювати, як події контрольованого оброблення уточнюють цей опис за відомих умов запуску. Інформаційна технологія виявлення аномальних кодів у фото- та відеооб’єктах вебресурсів повинна організовувати послідовність дій, зберігаючи підстави рішення та межі доступних свідчень.

Для експериментального оцінювання важливо перевіряти саме реалізований зв’язок між цими складниками. Результат окремої зорової гілки не характеризує структурний аналіз, а відсутність повного журналу подій не дає підстав стверджувати про перевірку всього мультимодального контуру. Порівняння має зберігати однакову одиницю розподілу, умови підготовки та правило прийняття рішення, щоб зміна показника відповідала зміні досліджуваного складника. Пояснення повинно відтворювати, які статичні та процесні свідчення вплинули на висновок і які дані були недоступні. Такий контракт не задає наперед бажаного результату, але робить твердження про модель, методи й технологію співвідносними з фактичним обсягом перевірки. На цій основі задача дослідження може бути сформульована без ототожнення очікуваних властивостей із уже встановленими експериментальними висновками.

Задача дослідження полягає у формуванні та операційному використанні мультимодального представлення прихованої загрози для виявлення аномальних кодів. Її розв’язання потребує формалізації маршруту об’єкта у вебсистемі, розроблення моделі мультимодального представлення прихованої загрози, удосконалення методу спектрально-зорового та структурного аналізу, формування методу поведінкового узгодження ознак різного походження і подальшого розвитку інформаційної технології їх послідовного застосування. Експериментальне оцінювання має охоплювати матрицю помилок, частку правильних рішень (Accuracy), точність позитивного виявлення (Precision), повноту (Recall), гармонійне поєднання точності й повноти ($F_1$), частки хибнопозитивних (FPR) і хибнонегативних (FNR) рішень, площу під ROC-кривою (ROC-AUC) за наявності неперервних оцінок та час оброблення. Сила висновків повинна відповідати фактичним даним; припущення не можуть замінювати відсутній експеримент або незалежну перевірку.

\section*{Висновки до розділу 1}
\addcontentsline{toc}{section}{Висновки до розділу 1}

Визначено причинний маршрут прихованого впливу: мультимедійний об’єкт надходить штатним каналом, містить доступну певному компоненту послідовність, проходить операцію її інтерпретації та лише після цього може створити спостережуваний прояв. Звідси випливає, що форматна коректність і зовнішня правдоподібність характеризують окремі етапи маршруту, але не доводять безпечності об’єкта для всіх компонентів вебсистеми. Для збереження свідчень потрібні незмінне початкове подання, зв’язок похідних копій із джерелом і реєстрація умов оброблення.

Систематизовано просторове, частотне, вейвлетне, часове, кодекове, структурне, службове й комбіноване розміщення за механізмом зміни, спостережуваними ознаками та умовами їх утрати. Установлено, що штатне перетворення змінює доступне подання: руйнування прихованої послідовності в одному каналі не виключає залишкового, структурного або подієво-поведінкового прояву в іншому. Для відео одиницею висновку має залишатися цілісний об’єкт зі збереженими часовими позиціями й доступними кодековими ознаками.

Порівняння статистичних, структурних, спектрально-зорових, відеокодекових, поведінкових і мультимодальних методів показало, що ці групи формують не взаємозамінні рішення, а взаємодоповнювальні джерела свідчень. Їхній результат залежить від джерела носіїв, способу вбудовування, кодека, доступних модальностей і протоколу перевірки. Критичними умовами доказовості визначено групово незалежний розподіл споріднених об’єктів, контроль зміщення джерела, явну маску доступності, локалізацію, калібрування робочої точки та вимірювання повного шляху оброблення.

Наукову прогалину визначено як недостатню опрацьованість узгодженого представлення внутрішньої будови об’єкта, слабких спектрально-зорових змін, службових і часових даних та реакції вебсистеми зі збереженням їхнього походження. Для її опрацювання обґрунтовано чотири взаємопов’язані складники: одну модель мультимодального представлення прихованої загрози, удосконалений метод спектрально-зорового та структурного аналізу, запропонований метод поведінкового узгодження ознак різного походження та одну інформаційну технологію їх послідовного застосування. Дослідницьке припущення пов’язує їх через спільний ідентифікатор, локалізацію, причинний контекст і маску доступності.

Визначено мінімальні свідчення для перевірки кожного складника й зіставлено їх із наявним корпусом. Доступні агреговані результати характеризують лише описану прототипну конфігурацію й дають змогу арифметично перевірити частину класифікаційних і часових показників. Вони не підтверджують повної навчуваної подієво-поведінкової гілки, наскрізного часу п’яти етапів або переносимості на незалежні дані. Тому подальші твердження про модель, два методи й інформаційну технологію мають залишатися в межах фактично відтворених свідчень.
